\section{随机变量变换与逆变换抽样}
\label{sec:05-Probability/03-Random-Variables/04}

写仿真的时候，标准库能给的只有一样东西：\([0,1]\) 上的均匀随机数。可你要的往往是别的——订单金额服从对数正态，故障间隔服从指数，测量噪声服从正态。中间那步换算由谁完成？

这个疑问属于一类更大的问题。已知 \(X\) 的分布，令 \(Y=g(X)\)，求 \(Y\) 的分布。摄氏度换华氏度是它，电压平方成功率是它，零件尺寸落不落在公差带里也是它（此时 \(g\) 只取 0 和 1）。开头那个采样需求则是同一件事反着提：先指定 \(Y\) 该服从什么分布，再回头找 \(g\)。

\subsection{兜底的办法只有一个}

不管 \(g\) 长什么样，退回分布函数总是安全的：
\[
F_Y(y)=P\bigl(g(X)\le y\bigr).
\]
右边是一个关于 \(X\) 的事件，把它解成 \(X\) 的不等式，再套上 \(X\) 已知的分布，就完事了。设 \(X\sim\mathcal N(0,1)\)、\(Y=X^2\)，对 \(y>0\) 有 \(X^2\le y\) 等价于 \(-\sqrt y\le X\le\sqrt y\)，于是 \(F_Y(y)=2\Phi(\sqrt y)-1\)。整个过程没有动用任何新定理，唯一的技术含量在解不等式。

\(g\) 严格递增且可逆时形状最干净：\(g(X)\le y\) 与 \(X\le g^{-1}(y)\) 同义，直接给出
\[
F_Y(y)=F_X\bigl(g^{-1}(y)\bigr).
\]
严格递减时不等号要翻，写成 \(F_Y(y)=P\bigl(X\ge g^{-1}(y)\bigr)\)。这个方向容易记反，每次从事件推一遍比背结论可靠。

离散情形连积分都省了：把所有映到同一个 \(y\) 的原像的质量加起来，\(P(Y=y)=\sum_{x:\,g(x)=y}P(X=x)\)。掷骰子取模 2 得到公平硬币，就是三个偶数并成一堆、三个奇数并成另一堆。

\subsection{导数因子是密度的伸缩比}

分布函数已经够用，但人更愿意看密度，因为密度看得见形状。对 \(F_Y\) 求导当然行，还有一条更说明问题的路：在小段上做质量守恒。

\(X\) 轴上一小段 \([x,x+dx]\) 承载的概率约为 \(f_X(x)\abs{dx}\)；变换把它推到 \(Y\) 轴上的 \([y,y+dy]\)，那里承载的概率约为 \(f_Y(y)\abs{dy}\)。两段装的是同一批随机结果，概率必须相等：
\[
f_Y(y)\abs{dy}\approx f_X(x)\abs{dx}.
\]
若 \(g\) 在 \(x\) 附近可微且 \(g'(x)\ne0\)，两段长度之比就是导数；把 \(x\) 换回 \(g^{-1}(y)\) 并整理，得
\[
f_Y(y)=f_X\bigl(g^{-1}(y)\bigr)\,\abs{\bigl(g^{-1}\bigr)'(y)}.
\]

三个条件各挡一件事：可微保证局部有比值可谈，导数非零保证 \(g\) 在该点附近一一对应、\(g^{-1}\) 确实存在，绝对值保证密度非负——递减变换的导数是负数，少了它密度会算成负的。

公式读一遍就知道两项分别在干什么。\(f_X\) 那一项回答``原来那个位置有多少概率''，导数那一项回答``那一小段被拉长还是压短了''。取 \(X\sim\Uniform(0,1)\)、\(Y=2X+3\)：区间从长度 1 拉成长度 2，密度就从 1 掉到 \(1/2\)，围出的面积仍是 1。同样的质量铺到两倍长的线上，线密度减半。

这个导数因子不是概率论的发明。定积分换元时冒出来的 \(\abs{dx/du}\)（见\hyperref[sec:02-Calculus/04-Integration/03]{``换元与分部积分：反向识别求导规则''}）、多重积分换坐标时冒出来的 Jacobian 行列式（见\hyperref[sec:02-Calculus/06-Multivariable-Calculus/07]{``多重积分与变量替换''}）、以及\hyperref[sec:03-Linear-Algebra/04-Determinants/01]{``行列式与体积：线性变换是放大、翻转，还是压扁空间''}里那个把体积缩放比压成一个数的行列式，说的是同一件事：一个变换把局部的``大小''改了多少倍。概率守恒，密度不守恒，差的正是这个倍数。上一节说密度带单位，也是这条式子的推论——把毫秒换成秒就是一次线性变换，导数因子 1000 全落在密度上。

多维情形照抄，只把导数换成 Jacobian 行列式的绝对值 \(\abs{\det J}\)。normalizing flow 把一串可逆变换叠起来构造复杂分布，训练时写进对数似然的那一项 \(\log\abs{\det J}\)，来源就在这里；那类模型的设计难点，多半是让这个行列式算得够快。

\subsection{一个 \texorpdfstring{\(y\)}{y} 有几个原像}

\(g\) 不是一一映射时，同一个 \(y\) 可能有好几个 \(x\) 与之对应，每一处都得交一份质量：
\[
f_Y(y)=\sum_{x:\,g(x)=y}\frac{f_X(x)}{\abs{g'(x)}}.
\]
做法是把定义域切成若干段，每段上 \(g\) 单调，分别套单调公式再相加。平方变换是最常见的例子——正负两侧的电压给出同样的功率，只代正根会丢掉一半概率质量。

切段时留意导数为零的点，它既是分段的界，也是密度可能爆掉的地方。\(Y=X^2\) 在 \(y\to0^{+}\) 处密度趋于无穷，因为 \(g'(x)=2x\) 在原点附近把一大片 \(x\) 挤进极短的一段 \(y\)。密度无界并不违规，它在那附近的积分仍然有限。

\subsection{只会掷均匀数的机器}

现在把箭头调头。手上只有 \(U\sim\Uniform(0,1)\)，想要服从指定 \(F\) 的样本，该找哪个 \(g\)？

上一节末尾已经把答案放在那里了：取 \(g=F^{-1}\)。验证只有一行，
\[
P\bigl(F^{-1}(U)\le x\bigr)=P\bigl(U\le F(x)\bigr)=F(x),
\]
第二个等号用的是均匀分布的定义 \(P(U\le u)=u\)，而 \(F(x)\) 总落在 \([0,1]\) 内，正好用得上。

\begin{center}
\begin{tikzpicture}[x=1.6cm,y=3.2cm,font=\footnotesize,>=stealth]
  \draw[->] (-0.06,0) -- (3.55,0) node[right] {\(x\)};
  \draw[->] (0,-0.03) -- (0,1.15) node[above] {\(u\)};
  \draw (-0.035,1) -- (0.035,1);
  \node[left,black!65] at (-0.05,1) {\(1\)};
  \foreach \u/\q in {0.1/0.105,0.2/0.223,0.3/0.357,0.4/0.511,0.5/0.693,0.6/0.916,0.7/1.204,0.8/1.609,0.9/2.303}
    {\draw[black!28] (0,\u) -- (\q,\u) -- (\q,0);
     \fill[orange!85!black] (0,\u) circle (1.5pt);
     \fill[blue!65!black] (\q,0) circle (1.5pt);}
  \draw[blue!70!black,thick] plot[domain=0:3.4,samples=80] (\x,{1-exp(-\x)});
  \node[right,black!70] at (2.05,0.42) {\(F(x)=1-e^{-x}\)};
  \node[black!70,align=center,font=\scriptsize] at (-0.58,0.55) {纵轴上\\等距的 \(u\)};
  \node[black!62,font=\scriptsize] at (0.42,-0.075) {挤密};
  \node[black!62,font=\scriptsize] at (1.95,-0.075) {拉开};
  \node[black!70] at (1.7,-0.185) {横轴上的落点 \(x=F^{-1}(u)\)};
\end{tikzpicture}

\emph{纵轴上九个等距的点代表均匀随机数可能落到的高度。每个点横着走到曲线，再垂直落到横轴，就得到一个服从 \(\Exp(1)\) 的样本。等距进去，不等距出来：\(F\) 陡的地方，一小段概率只对应很短的一段 \(x\)，落点被挤密；\(F\) 接近水平的尾部，同样一小段概率被摊到很长的区间上，落点稀疏。密度的高低，就是这张图上纵轴刻度被压缩或拉伸的倍数。}
\end{center}

拿 \(\Exp(2)\) 走一遍：解 \(u=1-e^{-2x}\) 得 \(x=-\tfrac12\log(1-u)\)。工程实现里常直接写 \(x=-\tfrac12\log u\)，因为 \(1-U\) 与 \(U\) 同分布，省一次减法。

也可以用上一小节的密度公式反过来核对。\(f_U\equiv1\)，而 \(g=F^{-1}\) 的反函数就是 \(F\)，代入变换公式得 \(f_X(x)=1\cdot\abs{F'(x)}=f(x)\)，正是目标密度。两条路给出同一个结论并不意外：逆变换抽样只是单调变换公式的一个特例，图上那种疏密变化，就是导数因子在起作用。

\subsection{这个办法管不到的地方}

\(F\) 不连续或不严格单调时，把 \(F^{-1}\) 换成上一节那个广义逆 \(Q(u)=\inf\set{x\given F(x)\ge u}\)，结论一字不改。离散分布正是这种情形，实现上从``解方程''变成``在累积概率表里查 \(u\) 落在哪一级台阶''。

真正的限制在别处。正态分布的 \(\Phi\) 没有初等闭式反函数，实践中或者用有理函数逼近 \(\Phi^{-1}\)，或者改用 Box--Muller 这类不碰反 CDF 的专用算法。更硬的一堵墙是维度：一维靠``把概率从小到大排开''建立的这套对应，在高维没有现成的对应物，何况目标分布常常只知道未归一化的形式。那时得换成在 Markov 链上采样（见\hyperref[sec:13-Machine-Learning/08-Sampling-and-Inference/02]{``MCMC 在无法独立采样时构造平稳链''}），与这里的确定性变换已经不是一条路线；至于拿样本去逼近期望的整套做法，见\hyperref[sec:13-Machine-Learning/08-Sampling-and-Inference/01]{``Monte Carlo 用随机样本逼近期望''}。

变换还有一个用法在深度学习里出现得很频繁。要对一个含参数的分布求梯度，直接对``采样''这个动作求导无从下手；写成 \(X=\mu+\sigma Z\)（其中 \(Z\) 与参数无关）之后，随机性被挪到了一个固定的源头上，参数只留在可微的变换里，梯度就能穿过去。这是变分自编码器的重参数化技巧，逆变换抽样是它最朴素的形态，\hyperref[ch:128]{``变分自编码器：把不可积的后验换成可优化的下界''}会展开这条线。

到这里，单个随机变量的分布已经能写、能变换、能采样。可一份完整的分布不便于比较，也不便于放进公式——报表上真正被拿来沟通的是几个数字：中心在哪、波动多大、两个量是否同向。下一章把分布再压一层。
